\FourSubsection{Rotated Periodic Boundary Condition}
\label{SEC:BC:Numerical:Periodic:Inputs:Rotated}

Rotated periodic boundary conditions are useful for problems containing
rotational symmetry.
%
This is particularly prevalent in turbomachinery flows where simulations of a
single periodic rotor blade can save significant computational expense compared
to simulations of the full 360-degree rotor wheel with dozens of rotor blades.
%
There are three subcomponent arrays of \BCinptext[i]{Periodic} that can be used
to specify the rotational transformation between a pair of rotated periodic
boundary conditions:
%
\begin{InputExample}
  \item[{\BCinpBC[nmlBC::BCInput/Periodic/RotationCenter][i]%&
                 {Periodic\%RotationCenter(1:3)}%
                 [Periodic\%RotationCenter]}]
  \item[{\BCinpBC[nmlBC::BCInput/Periodic/RotationAngle][i]%&
                 {Periodic\%RotationAngle(1:3)}%
                 [Periodic\%RotationAngle]}]
  \item[{\BCinpBC[nmlBC::BCInput/Periodic/RotationAxis][i]%&
                 {Periodic\%RotationAxis(1:4)}%
                 [Periodic\%RotationAxis]}]
\end{InputExample}
%
%One may use either \BCinptext[i]{Periodic\%RotationAngle} or
%\BCinptext[i]{Periodic\%RotationAxis} to define rotated periodic boundary
%conditions, but not both.
%
One may use either \inptext{RotationAngle} or \inptext{RotationAxis} to define
rotated periodic boundary conditions, but not both.
%
If both are specified, \inptext{RotationAxis} will be ignored and
\inptext{RotationAngle} will be used.
%
%For both methods, the input vector \BCinptext[i]{Periodic\%RotationCenter(1:3)}
%gives the rotation origin (in \key[key:GridUnits]{grid units}) used
%for defining the rotation angles between rotated periodic boundary conditions.
%
For both methods, the input vector \inptext{RotationCenter(1:3)} gives the
coordinates for the rotation origin (in \key[key:GridUnits]{grid units}) used
for defining the rotation angles between rotated periodic boundary conditions.

%%%%%%%%%%%%%%%%%%%%%%%%%%%%%%%%%%%%%%%%%%%%%%%%%%%%%%%%%%%%%%%%%%%%%%%%%%%%%%%%

The first way to specify the rotational transformation between a pair of rotated
periodic boundary conditions is through the input vector
\inptext{RotationAngle(1:3)}.
%\BCinptext[i]{Periodic\%RotationAngle(1:3)}.
%
This vector gives the rotation angles (in degrees) about each of the coordinate
axes from the current periodic boundary interface to the connecting periodic
boundary interface.
%
If rotating about more than one axis, the rotation is performed first about the
$x$-axis, then the $y$-axis, and finally the $z$-axis.
%
For a given periodic pair with rotational periodicity, the following relation
needs to be true:
%
\begin{InputExample}
  \item[\small{\BCinptext[i]{Periodic\%RotationAngle(:) = -}
               \BCinptext[j]{Periodic\%RotationAngle(:)}}]
\end{InputExample}

%%%%%%%%%%%%%%%%%%%%%%%%%%%%%%%%%%%%%%%%%%%%%%%%%%%%%%%%%%%%%%%%%%%%%%%%%%%%%%%%

The second way to specify the rotational transformation between a pair of
rotated periodic boundary conditions is through the input vector
\inptext{RotationAxis(1:4)}.
%\BCinptext[i]{Periodic\%RotationAxis(1:4)}.
%
The first 3 values of this array define a vector representing the axis of
rotation for this rotated periodic boundary condition, and the last value of the
array defines the rotation angle (in degrees) about this axis.
%
The magnitude of the vector defining the axis of rotation is not used; the sole
purpose of this vector is to define the direction of the rotation axis.
%
For a given periodic pair with rotational periodicity, one of the two following
sets of relations needs to be true:
%
\begin{InputExample}
  \item[\small{\BCinptext[i]{Periodic\%RotationAxis(1:3) = }
               \BCinptext[j]{Periodic\%RotationAxis(1:3)}}]
  \item[\small{\BCinptext[i]{Periodic\%RotationAxis(4) = -}
               \BCinptext[j]{Periodic\%RotationAxis(4)}}]
\end{InputExample}
%
or
%
\begin{InputExample}
  \item[\small{\BCinptext[i]{Periodic\%RotationAxis(1:3) = -}
               \BCinptext[j]{Periodic\%RotationAxis(1:3)}}]
  \item[\small{\BCinptext[i]{Periodic\%RotationAxis(4) = }
               \BCinptext[j]{Periodic\%RotationAxis(4)}}]
\end{InputExample}

%%%%%%%%%%%%%%%%%%%%%%%%%%%%%%%%%%%%%%%%%%%%%%%%%%%%%%%%%%%%%%%%%%%%%%%%%%%%%%%%

For example, say you have a rotor fan with 15 blades with the rotation axis
aligned with the Cartesian $x$-axis, and you would like to run a simulation of
just a single blade using rotational periodicity.
%
Using \inptext{RotationAngle}, your inputs for boundary condition $i$ would
look like
%
\begin{InputExample}
  \item[{\BCinptext[i]{Periodic\%RotationAngle(1:3) = [ 24.0, 0.0, 0.0]}}]
\end{InputExample}
%
and the inputs for boundary condition $j$ would look like
%
\begin{InputExample}
  \item[{\BCinptext[j]{Periodic\%RotationAngle(1:3) = [-24.0, 0.0, 0.0]}}]
\end{InputExample}
%
If you wanted to use \inptext{RotationAxis} instead, your inputs for boundary
condition $i$ would look like
%
\begin{InputExample}
  \item[{\BCinptext[i]{Periodic\%RotationAxis(1:4) = [ 1.0, 0.0, 0.0, 24.0]}}]
\end{InputExample}
%
and the inputs for boundary condition $j$ would look like
%
\begin{InputExample}
  \item[{\BCinptext[j]{Periodic\%RotationAxis(1:4) = [ 1.0, 0.0, 0.0,-24.0]}}]
\end{InputExample}

